\documentclass[11pt,a4paper]{article}
\usepackage[margin=2cm]{geometry}
\usepackage{booktabs,tabularx,xltabular,array,enumitem,xcolor,hyperref}
\usepackage[T1]{fontenc}
\usepackage{lmodern}
\hypersetup{colorlinks=true,linkcolor=blue!50!black,urlcolor=blue!50!black}
\newcolumntype{L}{>{\raggedright\arraybackslash}X}
\setlist{nosep,leftmargin=*}
\newcommand{\flag}[1]{\textcolor{red}{[#1]}}

\title{Microprocessor-Based Product Design 2026/27\\[2pt]\large Course plan: sessions and assignments}
\author{Yohai Ackerman (TA), draft for discussion with Prof.\ Izhak Bucher}
\date{Draft of \today}

\begin{document}
\maketitle

\section{Sessions}

\noindent
\begin{xltabular}{\textwidth}{@{}c l L L >{\raggedright\arraybackslash}p{1.6cm}@{}}
\toprule
\# & Lecture date & Lecture & Tutorial & Milestone \\
\midrule
\endhead
\bottomrule
\endfoot
1 & 2.11.2026 & Intro and product mindset; logic gates, flip-flops, counters, edge vs.\ level; integers in binary/hex, two's
    complement, overflow.
  & First STM32 application: CubeMX $\to$ build $\to$ flash, blink LD2, \texttt{printf} to PuTTY, breakpoints and
    live expressions.
  & \\
\addlinespace
2 & 9.11.2026 & Python $\to$ C: types (\texttt{stdint}), casting, scope, \texttt{\#define}/\texttt{const}, operators and
    precedence, shifts and masks, pointers, structs, registers as pointers.
  & Numbers and bits on the board: predict-then-run wrap-around, signed/unsigned in the debugger, binary printer, the
    four bit idioms (set/clear/toggle/test).
  & A1 out \\
\addlinespace
3 & 16.11.2026 & \textcolor{blue}{The IO board as a black box: RGB LEDs (3 chips $\times$ 8 channels), colour mixing,
    gamma, patterns as data;} buttons (active-low byte), edge detection, debouncing; ``Catch the Dot''. Timers and
    timing: \texttt{HAL\_Delay} vs.\ \texttt{HAL\_GetTick}, non-blocking multi-rate timing. The UART protocol: baud
    rate, the 8N1 frame (start bit, 8 data bits, stop bit), the time to send one character.
    \flag{Yohai to discuss with Bucher: move timers, \texttt{HAL\_GetTick} and the UART protocol into lecture 3, since
    Assignment 1 needs them; the blue part could shrink to make room.}
  & RGB and buttons: \texttt{set\_rgb} helper, walking dot with masks, saturating brightness up/down, one-press-one-action.
  & \\
\addlinespace
4 & 23.11.2026 & Timers: counter, prescaler, period, polling vs.\ interrupts, priorities, \texttt{volatile}; software state machines;
    \texttt{HAL\_GetTick} non-blocking timing.
  & Multi-rate blink from a timer interrupt (no \texttt{HAL\_Delay}); state-machine skeleton (\texttt{enum} +
    \texttt{switch}).
  & Quiz 1 (first hour) \\
\addlinespace
5 & 30.11.2026 & The I2C bus opened: protocol, addressing, ACK/NACK, pull-ups; HAL I2C; register-level drivers for PCA9634,
    PCA9538A, PCA9685 (servo pulses).
  & I2C bus scan: detect the five IO-board chips (3$\times$ PCA9634 RGB LED drivers, PCA9538A buttons, PCA9685 servo driver); rewrite \texttt{set\_rgb} on raw \texttt{HAL\_I2C\_Mem\_Write}; servo sweep.
  & A1 due \\
\addlinespace
6 & 7.12.2026 & PWM and input capture: duty cycle, H-bridge preview; SR04 trigger/echo by input capture; EXTI on the
    optical sensor (edge triggering revisited).
  & SR04 single distance measurement; EXTI interrupt with a counter; PWM on one motor channel.
  & A2 out \\
\addlinespace
7 & 14.12.2026 & DC motors and encoders: H-bridge drive, quadrature decoding (encoder timer mode), speed measurement, recording
    a step response at a fixed rate.
  & One wheel in the air: open-loop PWM, encoder to rad/s, record and export to MATLAB.
  & \\
\addlinespace
8 & 21.12.2026 & PI speed control: model fitting (V $\to$ rad/s), design, discretisation and sample rate, PI inside a 100\,Hz
    interrupt, saturation and anti-windup, STM32CubeMonitor.
  & Control workflow: fit, \texttt{sisotool}, \texttt{c2d}, implement the difference equation in C,
    compare simulated vs.\ measured step.
  & Quiz 2 (first hour) \\
\addlinespace
9 & 28.12.2026 & UART command protocol: interrupt-driven receive, command interpreter, ASCII vs.\ binary framing; MATLAB
    \texttt{serialport} and an App Designer GUI.
  & UART command parser; minimal MATLAB GUI that sends commands and plots telemetry.
  & A2 due \\
\addlinespace
10 & 4.1.2027 & Mecanum kinematics: inverse and forward kinematics, odometry; trajectory generation (trapezoidal profiles, waypoints,
     piecewise-constant $(v_x,v_y,\omega)$ with output-compare timing).
   & Kinematics exercises; driving the cart on the floor; a trajectory array generated in MATLAB and sent to the cart.
   & \\
\addlinespace
11 & 11.1.2027 & Simulink to C: discrete blocks and sample rates, exporting a controller or filter, integrating it in CubeIDE;
     digital filters (median, IIR).
   & Export a Simulink PI/filter to C and run it on the cart.
   & A3 out \\
\addlinespace
12 & 18.1.2027 & System integration and case studies: scanned sonar map, obstacle response, wrap-up; final hints for A3.
   & Open project clinic.
   & \\
\addlinespace
13 & 25.1.2027 & Lab meeting: open questions.
   &
   & A3 due 3.3.2027 \\
\end{xltabular}

\section{Grading}

\noindent
\begin{tabularx}{\textwidth}{@{}l L r@{}}
\toprule
Component & Contents & Weight \\
\midrule
Quiz 1 & Sessions 1 to 3: logic and numbers, C types and casting, bit operations, IO board LEDs and buttons & 10\% \\
Quiz 2 & Sessions 4 to 7: timers and interrupts, state machines, I2C, PWM and input capture, motors and encoders & 10\% \\
Assignment 1 & Variable types, wrap-around, bit operations and shifting, logic, LED blinking, RGB LEDs & 20\% \\
Assignment 2 & DC motors and PWM, encoders, timer interrupts, PI speed control, UART (wheels in the air) & 25\% \\
Assignment 3 & Driving on the ground, ultrasonic sensing, MATLAB GUI, Simulink to C & 35\% \\
\bottomrule
\end{tabularx}

\section{Assignments}

\subsection{General instructions}
\begin{itemize}
  \item \textbf{Report format:} PDF only, digitally typed (not handwritten). The report may be written in Hebrew or English.
  \item \textbf{Report structure:}
  \begin{itemize}
    \item Title page: assignment number, names of the submitters, ID numbers, emails and date.
        \item An answer to each part. Do not Copy full codes section into the report, but DO explain what have you done (flow chart \ pseudo code). 
  \end{itemize}
  \item \textbf{Code:} Use function for the code to be organized -  each function must have a documentaion in the beginning - what it is suppose to do. \\Document and explain your code, after a line of a code. \\ Write in the user-defined section, as taught.\\ Be oragnized. Unused codes section should not be included!
  \item \textbf{Submission:} a single ZIP file named \texttt{lab\#\_studentname1\_studentname2.zip}, containing the code
        folders of the different parts, the report and a video demonstrating the system in operation.
  \item The grade is based, among other things, on the efficiency and clarity of the code.
  \item An engineering approach is expected: if you need to add a task, a code section, a state or a detail in the report
        in order to fulfil the requirements, do it. Creativity and elegant writing will be rewarded.
\end{itemize}
\subsection{Assignment 1: ``No wheels''}

\textbf{Window:} out at session~2, due at session~5.\quad
\textbf{Hardware:} NUCLEO-G474RE and the I2C IO board only. The motors are not powered.\quad
\textbf{Allowed tools:} \texttt{HAL\_Delay}/\texttt{HAL\_GetTick}, GPIO, \texttt{printf} over LPUART1, the debugger, and
the provided IO-board functions. The I2C internals stay a black box until session~5.

\medskip\noindent
\textbf{Provided code.} The RGB LEDs and the buttons sit on the I2C bus, which is taught only in session~5. Until then,
use the provided files \texttt{io\_board.h} and \texttt{io\_board.c}: copy \texttt{io\_board.h} into \texttt{Core/Inc}
and \texttt{io\_board.c} into \texttt{Core/Src} (in CubeIDE, refresh the project with F5 afterwards). In
\texttt{main.c}, add \texttt{\#include "io\_board.h"} inside \texttt{USER CODE BEGIN Includes} (include the
\texttt{.h} file only, never the \texttt{.c} file), and call \texttt{ioBoardInit()} once inside
\texttt{USER CODE BEGIN 2}. CubeIDE compiles \texttt{io\_board.c} automatically. To make \texttt{printf} print to PuTTY, copy the function in
Appendix~A1 into your \texttt{main.c}.

\medskip\noindent
\textbf{Learning outcome.} Students work confidently in C on real hardware. They choose variable types deliberately,
predict and explain overflow and wrap-around, manipulate individual bits with shifts and masks, combine inputs with
logic, and produce timed LED behaviour on the on-board LED and the RGB LEDs.

\paragraph{Part A: Variable types and wrap-around.}
\begin{enumerate}[label=A.\arabic*, leftmargin=2.2em]
  \item Increment a \texttt{uint8\_t}, an \texttt{int8\_t} and a \texttt{uint16\_t} past their limits. Watch them in Live
        Expressions. Print them to PuTTY in decimal, hex and binary, and capture it in a video (think how to do it so
        that the viewer can clearly see the wrap-around). In your report, describe what happens when a variable exceeds
        its limit, and explain why.
  \item Explain how to interpret a hex number as a decimal value, and the difference between the signed and the
        unsigned cases. Show examples for \texttt{0x7F} (\texttt{uint8\_t} and \texttt{int8\_t}), \texttt{0xF0}
        (\texttt{uint8\_t} and \texttt{int8\_t}) and \texttt{0xFFF6} (\texttt{uint16\_t} and \texttt{int16\_t}).
        In the code, assign each value in its hex form to variables of both types, and print them to PuTTY in hex,
        binary and decimal. Explain what you see when a negative value is printed in hex.
  \item Integer division and casting. Two students took the MP course. One attended the classes and got a good grade,
        97. The other did not, and got a bad grade, 58. Store each grade in an unsigned integer variable, then compute
        their average: first in a way that gives a wrong result because of integer arithmetic, then in a way that gives
        the correct result. Show both in the code, and explain the difference in your report.
  % \item A brightness value that wraps around when it is increased past its maximum, then fixed with saturation.
\end{enumerate}

\paragraph{Part B: Bit operations and shifting.}
\begin{enumerate}[label=B.\arabic*, leftmargin=2.2em]
  \item Write the four bit idioms as functions (not macros). Each takes a value and a bit number, and returns the
        result: \texttt{bitSet}, \texttt{bitClear} and \texttt{bitToggle} return the modified value, and
        \texttt{bitTest} returns 0 or 1. For example: \texttt{GPIOA->ODR = bitSet(GPIOA->ODR, 5);}
        Use them to operate the Nucleo's green LED (LD2, bit 5 of \texttt{GPIOA->ODR}) directly through the register,
        without HAL GPIO functions, in this sequence: 1\,s of rest, SET; 1\,s of rest, CLEAR; then 3 times (1\,s of
        rest, TOGGLE); then TEST. Show the test result on LED 0 of the IO board: lit when the result is 1, off when it
        is 0.
  \item Walking single LED. Keep the state of the 8 RGB LEDs in one \texttt{uint8\_t} variable, in which bit $n$
        represents LED $n$ (1 = lit). Light one LED at a time, with 1\,s between transitions: LED 0, then LED 1, and so
        on up to LED 7, then back to LED 0, forever. Print the variable in binary at every step.
  \item Walking pair of LEDs. The same, with two adjacent LEDs lit: LEDs 0 and 1, then 1 and 2, and so on; after
        LEDs 6 and 7 come LEDs 7 and 0, then 0 and 1 again.
\end{enumerate}
\smallskip\noindent Rules for B.2 and B.3: the variable may change only through bit operations (shifts, OR, AND). Choosing the
pattern with \texttt{if}/\texttt{switch} chains, counters mapped to LEDs, or tables of precomputed patterns is not
accepted. The LEDs are then lit from the variable, by testing its bits. In your report, explain how your code moves
the pattern past bit 7 without losing it.
\smallskip
\begin{enumerate}[label=B.\arabic*, leftmargin=2.2em, start=4]
  \item Read the raw button byte with \texttt{readButtonsRaw()}. Using bit operations, convert it into a value in
        which bit $n$ is 1 while button $n$ is held: invert the active-low bits, then clear bits 4 to 7 with a mask.
        Count how many buttons are held (0 to 4) by testing each bit, and print the count whenever it changes.
        \flag{For Bucher and Yohai, delete before publishing: \texttt{io\_board.h} also provides
        \texttt{readButtonsPressed()}, which already returns \texttt{(\textasciitilde raw) \& 0x0F} and so solves this
        item. Consider providing a copy of \texttt{io\_board} without it.}
        \flag{Also: \texttt{setLedMask()} in \texttt{io\_board.h} lights LEDs straight from a bit mask, which shortens
        the display step of B.2 and B.3.}
\end{enumerate}

\paragraph{Part C: Input triggering.}
\begin{enumerate}[label=C.\arabic*, leftmargin=2.2em]
  \item Count presses of button 0 and show the count on the 8 RGB LEDs as an 8-bit binary number (bit $n$ = LED
        $n$); print it as well. Do it three times:
        \begin{enumerate}[label=(\alph*)]
          \item Level: increase the counter on every pass of the main loop while button 0 is held. Describe what
                happens.
          \item Edge: increase the counter only when button 0 changes from released to pressed
                (\texttt{now \& \textasciitilde before}), polling as fast as possible, with no delay in the loop.
                Press the button exactly 20 times and record the count. Repeat 3 times.
          \item Debounced edge: accept a change of the button state only after it has been stable for 20\,ms. Press
                exactly 20 times, 3 times over.
        \end{enumerate}
        In your report, compare the three results and explain each difference.
        \flag{For Bucher and Yohai, delete before publishing: (b) shows bouncing only if the IO-board buttons
        bounce and no hardware filter hides it. Verify on the cart with state 17 of the hardware test before
        publishing; if the I2C buttons show no bounce, use the Nucleo button B1 with an EXTI interrupt instead.}
\end{enumerate}

\paragraph{Part D: LED blinking and timing.}
\begin{enumerate}[label=D.\arabic*, leftmargin=2.2em]
  \item Write a loop that toggles LD2, prints a status line of about 60 characters to PuTTY with \texttt{printf}, and
        then waits with \texttt{HAL\_Delay(500)}; stop after 200 toggles. Before running it, calculate how long the
        200 toggles will take; then run it and measure. In your report, explain why the total is not exactly
        200 $\times$ 500\,ms, using the UART frame format and the baud rate.
  \item Blink LD2 (PA5) at a rate set by a global variable that can be changed in the debugger.
  \item Two LEDs blinking at different rates, done twice: first with \texttt{HAL\_Delay}, then non-blocking with
        \texttt{HAL\_GetTick}. Measure and explain the difference.
\end{enumerate}
\paragraph{Part E: Integration and logic.}
Write a program with an \emph{armed} state that button 0 toggles (one press = one action, as in Part C). The trigger
condition is: the system is armed \textbf{and} the IR sensor (PA7) detects an obstacle.
\begin{itemize}
  \item While the condition holds, all 8 LEDs (white) grow brighter step by step, from off to full brightness, over
        about 2\,s. When full brightness is reached, the LEDs switch off, and stay off until the condition has become
        false and true again.
  \item If the condition stops holding during the brightening, the LEDs dim step by step from their current
        brightness until they are off.
\end{itemize}
Set the brightness of all LEDs with \texttt{setRGB(v, v, v)} from \texttt{io\_board.h} ($v$ = 0 to 255). Pace the steps
with \texttt{HAL\_GetTick}, so the inputs are read on every pass of the loop. In your report, give the truth table of
the trigger condition, and draw the behaviour as a state diagram. Creativity bonus applies here.

\subsubsection*{Assignment 1: project configuration (\texttt{.ioc})}
Everything needed to build the STM32CubeMX configuration for Assignment~1. The wheels are not used.

\paragraph{Creating the project.}
\begin{enumerate}
  \item STM32CubeIDE: \emph{File $\to$ New $\to$ STM32 Project}, tab \emph{Board Selector}, choose \textbf{NUCLEO-G474RE}.
  \item Answer \textbf{Yes} to ``Initialize all peripherals with their default mode''. This already configures LD2,
        the B1 button, LPUART1 and the debug pins.
  \item Apply the settings below, then \emph{Project $\to$ Generate Code}.
\end{enumerate}

\paragraph{System.}\leavevmode\\[2pt]
\begin{tabularx}{\textwidth}{@{}l L@{}}
\toprule
Item & Setting \\
\midrule
MCU & STM32G474RET6 (Cortex-M4F with FPU) \\
System clock & SYSCLK = 170\,MHz from the PLL (HSE source). HCLK = APB1 = APB2 = 170\,MHz \\
Debug & Serial Wire: PA13 (SWDIO), PA14 (SWCLK), PB3 (SWO). \textbf{Never reassign these pins.} \\
Time base & SysTick, 1\,ms (used by \texttt{HAL\_Delay} and \texttt{HAL\_GetTick}) \\
Board power & Power jumper on \textbf{5V\_STLK} (USB power) \\
\bottomrule
\end{tabularx}

\paragraph{Pins.}\leavevmode\\[2pt]
\begin{tabularx}{\textwidth}{@{}l l l L@{}}
\toprule
Pin & CubeMX signal & Label & Function \\
\midrule
PA5 & GPIO\_Output & LD2 & Green user LED on the Nucleo board \\
PC13 & GPIO\_EXTI13 & B1 & Blue user button on the Nucleo board (EXTI, rising edge) \\
PA2 & LPUART1\_TX & & Serial console to the PC (ST-LINK virtual COM port) \\
PA3 & LPUART1\_RX & & Serial console from the PC \\
PB8 & I2C1\_SCL & & I2C clock to the IO board \\
PB9 & I2C1\_SDA & & I2C data to the IO board \\
PA7 & GPIO\_Input, no pull & IR\_SENSOR & IR obstacle sensor at the front of the cart. Active-low: LOW = obstacle \\
\bottomrule
\end{tabularx}

\paragraph{Peripherals.}\leavevmode\\[2pt]
\begin{tabularx}{\textwidth}{@{}l >{\raggedright\arraybackslash}p{4.2cm} L@{}}
\toprule
Peripheral & Mode & Settings \\
\midrule
I2C1 & I2C & Standard mode, 100\,kHz \\
LPUART1 & Asynchronous & 115\,200 baud, 8 data bits, no parity, 1 stop bit \\
NVIC & & EXTI line [15:10] interrupt enabled (B1) \\
\bottomrule
\end{tabularx}

\paragraph{Devices on the I2C bus (IO board).}\leavevmode\\[2pt]
\begin{tabularx}{\textwidth}{@{}l l L@{}}
\toprule
Address (7-bit) & Chip & Role \\
\midrule
0x01 & PCA9634 & Red channel of LEDs 0 to 7 (chip channel $n$ = LED $n$) \\
0x02 & PCA9634 & Green channel of LEDs 0 to 7 \\
0x07 & PCA9634 & Blue channel of LEDs 0 to 7 \\
0x71 & PCA9538A & Four push buttons on bits 0 to 3, active-low \\
\bottomrule
\end{tabularx}

\paragraph{Serial terminal.} Connect PuTTY to the ST-LINK COM port at 115\,200 baud, 8 data bits, no parity, 1 stop
bit.

\subsubsection*{Appendix A1: code to copy}
\paragraph{printf to PuTTY.} Copy into \texttt{main.c}: the include into \texttt{USER CODE BEGIN Includes}, the function
into \texttt{USER CODE BEGIN 0}.
\begin{verbatim}
#include <stdio.h>

/* route printf to LPUART1 (PuTTY) */
int __io_putchar(int ch)
{
    HAL_UART_Transmit(&hlpuart1, (uint8_t *)&ch, 1, HAL_MAX_DELAY);
    return ch;
}
\end{verbatim}
End every printed line with \texttt{\textbackslash r\textbackslash n}: the output appears only when a line is complete.

\subsection{Assignment 2: ``Wheels in the air''}
To be defined.

\subsubsection*{Assignment 2: project configuration (\texttt{.ioc})}
The complete Assignment~1 configuration, plus the wheels (motors and encoders), the ultrasonic sensor and the servo
that turns it. Create the project as in Assignment~1. All timer clocks are 170\,MHz.

\paragraph{Pins from Assignment 1.} PA5 (LD2), PC13 (B1), PA2/PA3 (LPUART1), PB8/PB9 (I2C1), PA7 (IR\_SENSOR), with
the same settings as in Assignment~1.

\paragraph{Wheels: motors and encoders.} Each motor is driven by one PWM channel (speed) and one direction pin. Each
encoder is read by its own timer in encoder mode.\leavevmode\\[2pt]
\begin{tabularx}{\textwidth}{@{}c l l l L@{}}
\toprule
Motor & PWM pin & DIR pin & Encoder timer & Encoder pins (A, B) \\
\midrule
1 & PA0 = TIM2\_CH1 & PC2 (DIR1) & TIM1 & PC0 = TIM1\_CH1, PC1 = TIM1\_CH2 \\
2 & PA1 = TIM2\_CH2 & PB4 (DIR2) & TIM8 & PC6 = TIM8\_CH1, PC7 = TIM8\_CH2 \\
3 & PB10 = TIM2\_CH3 & PB11 (DIR3) & TIM3 & PA6 = TIM3\_CH1, PA4 = TIM3\_CH2 \\
4 & PA10 = TIM2\_CH4 & PB13 (DIR4) & TIM4 & PA11 = TIM4\_CH1, PA12 = TIM4\_CH2 \\
\bottomrule
\end{tabularx}

\smallskip\noindent
The DIR pins are GPIO\_Output, push-pull, no pull, initial level LOW. The level that drives a motor ``forward'' is SET
for motors 1 and 4 and RESET for motors 2 and 3, because the motors on the two sides of the cart are mirrored.

\paragraph{Ultrasonic sensor (HC-SR04).}\leavevmode\\[2pt]
\begin{tabularx}{\textwidth}{@{}l l L@{}}
\toprule
Pin & CubeMX signal & Function \\
\midrule
PB2 & TIM5\_CH1, PWM Generation CH1 & TRIG: a 10\,$\mu$s pulse every 100\,ms, generated by the timer \\
PC12 & TIM5\_CH2, Input Capture direct mode & ECHO: both edges are captured; the pulse width is the round-trip time \\
\bottomrule
\end{tabularx}

\smallskip\noindent
Distance: $d\,[\mathrm{mm}] = t_\mathrm{echo}\,[\mu\mathrm{s}] \times 343 / 2000$.

\paragraph{Peripherals.}\leavevmode\\[2pt]
\begin{tabularx}{\textwidth}{@{}l >{\raggedright\arraybackslash}p{4.2cm} L@{}}
\toprule
Peripheral & Mode & Settings \\
\midrule
I2C1 & I2C & Standard mode, 100\,kHz \\
LPUART1 & Asynchronous & 115\,200 baud, 8N1 \\
TIM2 & PWM Generation CH1 to CH4 & Prescaler 0, period 9\,999: 17\,kHz PWM, duty = CCR / 10\,000. Pulse 0, polarity
  high \\
\addlinespace
TIM1, TIM8, TIM3, TIM4 & Combined channels: Encoder Mode TI1 & Prescaler 0, period 65\,535: 16-bit up/down counter.
  Input polarity rising, filter 0 \\
\addlinespace
TIM5 & CH1 PWM Generation, CH2 Input Capture direct mode & Prescaler 9, period 1\,699\,999: tick = 1/17\,$\mu$s,
  period 100\,ms. CH1 pulse = 170 (10\,$\mu$s). CH2 polarity: both edges \\
\addlinespace
NVIC & & EXTI line [15:10] (B1), TIM2 and TIM5 global interrupts enabled \\
\bottomrule
\end{tabularx}

\paragraph{Devices on the I2C bus.} As in Assignment~1, plus:\leavevmode\\[2pt]
\begin{tabularx}{\textwidth}{@{}l l L@{}}
\toprule
Address (7-bit) & Chip & Role \\
\midrule
0x44 & PCA9685 & Servo pulse generator, 50\,Hz (PRESCALE = 121). Drives the servo that turns the ultrasonic sensor.
  Pulse width 900\,$\mu$s (one end), 1500\,$\mu$s (centre), 2100\,$\mu$s (other end) \\
\bottomrule
\end{tabularx}

\subsection{Assignment 3: ``Wheels on the ground''}
To be defined.

\subsection{Quizzes}
Two quizzes, each 10\% of the final grade, held in the first hour of the lecture: Quiz 1 in session 4 (23.11.2026), Quiz 2 in session 8 (21.12.2026).

\end{document}
